% defer/seqlock.tex
% mainfile: ../perfbook.tex
% SPDX-License-Identifier: CC-BY-SA-3.0

\section{顺序锁（sequence lock）}
\label{sec:defer:Sequence Locks}
%
\epigraph{就好像从头再来一样。}{John Lennon}

已发表的 sequence-lock
记录~\cite{10.1145/800212.806505,10.1145/359863.359878}
可以追溯到与读写锁（reader-writer lock）一样远，但 sequence locks 仍然相对鲜为人知。
Sequence locks 在 Linux 内核中用于必须被读者看到一致状态的
读多数据。然而，与读写锁不同，读者不排斥写者。
相反，像 hazard pointers 一样，sequence locks 迫使读者
在检测到并发写者的活动时\emph{重试}操作。
从\cref{fig:defer:Reader And Uncooperative Sequence Lock}可以看出，
设计使用 sequence locks 的代码时，让读者很少需要重试是很重要的。

\begin{figure}
\centering
\resizebox{3in}{!}{\includegraphics{cartoons/r-2014-Start-over}}
\caption{读者与不合作的 Sequence Lock}
\label{fig:defer:Reader And Uncooperative Sequence Lock}
\end{figure}

\QuickQuiz{
	为什么这个 sequence-lock 讨论不在\cref{chp:Locking}中，
	你知道的，关于\emph{锁}的那一章？
}\QuickQuizAnswer{
	sequence-lock 机制实际上是两种独立同步机制的组合，
	即序列计数和锁。
	事实上，序列计数机制在 Linux 内核中通过
	\co{write_seqcount_begin()} 和 \co{write_seqcount_end()}
	原语可以单独使用。

	然而，组合的 \co{write_seqlock()} 和
	\co{write_sequnlock()} 原语在 Linux 内核中使用得更多。
	更重要的是，如果你说``sequence lock''，
	比你说``sequence count''会有更多人理解你的意思。

	所以这一节标题为``Sequence Locks''以便人们仅从标题就能
	理解它讲的是什么，
	它出现在``延迟处理''中是因为 (1) 强调了``sequence locks''
	的``sequence count''方面，(2) 因为``sequence lock''远不仅仅是一个锁。
}\QuickQuizEnd

\begin{listing}
\begin{VerbatimL}
do {
	seq = read_seqbegin(&test_seqlock);
	/* read-side access. */
} while (read_seqretry(&test_seqlock, seq));
\end{VerbatimL}
\caption{Sequence-Locking 读者}
\label{lst:defer:Sequence-Locking Reader}
\end{listing}

\begin{listing}
\begin{VerbatimL}
write_seqlock(&test_seqlock);
/* Update */
write_sequnlock(&test_seqlock);
\end{VerbatimL}
\caption{Sequence-Locking 写者}
\label{lst:defer:Sequence-Locking Writer}
\end{listing}

sequence locking 的关键组件是序列号，
在没有更新者时具有偶数值，如果有更新正在进行则为奇数值。
读者可以在每次访问前后快照该值。
如果任一快照具有奇数值，或者两个快照不同，
则发生了并发更新，读者必须丢弃访问结果然后重试。
因此，读者使用 \co{read_seqbegin()} 和 \co{read_seqretry()}
函数（如\cref{lst:defer:Sequence-Locking Reader}所示）
来访问由 sequence lock 保护的数据。
写者必须在每次更新前后递增该值，
且同一时间只允许一个写者。
因此，写者使用 \co{write_seqlock()} 和 \co{write_sequnlock()}
函数（如\cref{lst:defer:Sequence-Locking Writer}所示）
来更新由 sequence lock 保护的数据。

结果是，受 sequence-lock 保护的数据可以有任意多个并发读者，
但同一时间只能有一个写者。
Sequence locking 在 Linux 内核中用于保护时间保持所用的校准量。
它还在路径名遍历中用于检测并发重命名操作。

\begin{listing}
\input{CodeSamples/defer/seqlock@impl.fcv}
\caption{Sequence-Locking 实现}
\label{lst:defer:Sequence-Locking Implementation}
\end{listing}

\cref{lst:defer:Sequence-Locking Implementation}中
展示了 sequence locks 的一个简单实现
（\path{seqlock.h}）。
\begin{fcvref}[ln:defer:seqlock:impl:typedef]
\co{seqlock_t} 数据结构在\clnrefrange{b}{e}中展示，
包含序列号和一个用于序列化写者的锁。
\end{fcvref}
\begin{fcvref}[ln:defer:seqlock:impl:init]
\Clnrefrange{b}{e}展示了 \co{seqlock_init()}，顾名思义，
它初始化一个 \co{seqlock_t}。
\end{fcvref}

\begin{fcvref}[ln:defer:seqlock:impl:read_seqbegin]
\Clnrefrange{b}{e}展示了 \co{read_seqbegin()}，
它开始一个 sequence-lock
\IXh{read-side}{critical section}。
\Clnref{fetch}获取序列计数器的快照，
\clnref{mb}将此快照操作排序在调用者的临界区之前。
最后，\clnref{ret}返回快照的值（最低有效位被清除），
调用者将把它传递给后续对 \co{read_seqretry()} 的调用。
\end{fcvref}

\QuickQuiz{
	为什么不让\cref{lst:defer:Sequence-Locking Implementation}中的
	\co{read_seqbegin()} 检查序列号值是否为奇数，
	如果是，则在内部重试而不是进入注定失败的
	读侧临界区？
}\QuickQuizAnswer{
	这将是一个合法的实现。

	但请记住
	\begin{enumerate*}[(1)]
	\item	这个额外检查是一个相对昂贵的条件分支，
	\item	它不能替代 \tco{read_seqretry()} 后面的检查，
		后者必须在临界区完成后发生，而且
	\item	Sequence locking 是为读多工作负载设计的，
		这意味着这个额外检查会减慢常见情况。
	\end{enumerate*}

	另一方面，在一个具有足够大比例更新和足够高开销读者的
	另一个世界中，在 \co{read_seqbegin()} 内部进行此检查
	可能更可取。

	\begin{fcvref}[ln:defer:seqlock:impl]
	当然，\cref{lst:defer:Sequence-Locking Implementation}的
	\clnref{read_seqbegin:mb,read_seqretry:mb}上的完整 memory barriers
	作为指令来说是相当重量级的，这表明添加检查的开销可能
	微不足道。
	只是在用户空间代码中，\co{membarrier()} 系统
	调用~\cite{Corbet2010membarrier,MathieuDesnoyers2017membarrier,Linuxmanpage2018sys-membarrier}
	可以用来消除读侧的 \co{smp_mb()} 开销，
	代价是更新侧 \co{membarrier()} 的额外开销。
	此功能正以``非对称屏障''的名称进入 C++ 标准~\cite{DavidGoldblatt2022asymmetricFences}。
	无论如何，这个技巧消除了读侧 \co{smp_mb()} 开销，
	这反过来使消除检查更有吸引力，
	再次回到\cref{lst:defer:Sequence-Locked Pre-BSD Routing Table Lookup}中
	展示的代码形式。
	\end{fcvref}

	同样的技巧也可以应用于 Linux 内核代码，使用
	\co{smp_call_function()} 等工具，至少在 Linux 内核的
	非实时构建中。
}\QuickQuizEnd

\begin{fcvref}[ln:defer:seqlock:impl:read_seqretry]
\Clnrefrange{b}{e}展示了 \co{read_seqretry()}，
如果自对应的 \co{read_seqbegin()} 调用以来至少有一个写者，
则返回 \co{true}。
\Clnref{mb}将调用者先前的临界区排序在\clnref{fetch}获取
序列计数器的新快照之前。
\Clnref{ret}检查序列计数器是否已更改，
换句话说，是否至少有一个写者，如果是则返回 \co{true}。
\end{fcvref}

\QuickQuizSeries{%
\QuickQuizB{
	为什么需要\cref{lst:defer:Sequence-Locking Implementation}
	\clnrefr{ln:defer:seqlock:impl:read_seqretry:mb}上的 \co{smp_mb()}？
}\QuickQuizAnswerB{
	如果省略它，编译器和 CPU 都有权将 \co{read_seqretry()} 调用
	之前的临界区移到该函数之下。
	这将阻止 sequence lock 保护临界区。
	\co{smp_mb()} 原语防止此类重排序。
}\QuickQuizEndB
%
\QuickQuizM{
	\cref{lst:defer:Sequence-Locking Implementation}中的代码
	能否使用更弱的 memory barriers？
}\QuickQuizAnswerM{
	在较旧版本的 Linux 内核中，不能。

	\begin{fcvref}[ln:defer:seqlock:impl]
	在非常新的 Linux 内核版本中，
	\clnref{read_seqbegin:fetch}可以使用
	\co{smp_load_acquire()} 代替 \co{READ_ONCE()}，
	这反过来允许删除\clnref{read_seqbegin:mb}上的 \co{smp_mb()}。
	类似地，\clnref{write_sequnlock:inc}可以使用
	\co{smp_store_release()}，例如如下：

\begin{VerbatimU}
smp_store_release(&slp->seq, READ_ONCE(slp->seq) + 1);
\end{VerbatimU}

	这将允许删除\clnref{write_sequnlock:mb}上的 \co{smp_mb()}。
	\end{fcvref}
}\QuickQuizEndM
%
\QuickQuizE{
	什么阻止了 sequence-locking 更新者饿死读者？
}\QuickQuizAnswerE{
	什么都没有。
	这是 sequence locking 的弱点之一，
	因此你应该只在读多情况下使用 sequence locking。
	当然，除非读侧饥饿（starvation）在你的场景中是可以接受的，
	在这种情况下，尽情使用 sequence-locking 更新吧！
}\QuickQuizEndE
}

\begin{fcvref}[ln:defer:seqlock:impl:write_seqlock]
\Clnrefrange{b}{e}展示了 \co{write_seqlock()}，
它简单地获取锁，递增序列号，
并执行 memory barrier 以确保此递增排序在调用者的临界区之前。
\end{fcvref}
\begin{fcvref}[ln:defer:seqlock:impl:write_sequnlock]
\Clnrefrange{b}{e}展示了 \co{write_sequnlock()}，
它执行 memory barrier 以确保调用者的临界区排序在
\clnref{inc}上序列号的递增之前，然后释放锁。
\end{fcvref}

\QuickQuizSeries{%
\QuickQuizB{
	如果有其他东西序列化写者，那么锁就不需要了，
	对吗？
}\QuickQuizAnswerB{
	在这种情况下，\co{->lock} 字段可以省略，
	就像 Linux 内核中的 \co{seqcount_t} 那样。
}\QuickQuizEndB
%
\QuickQuizE{
	为什么\cref{lst:defer:Sequence-Locking Implementation}
	\clnrefr{ln:defer:seqlock:impl:typedef:seq}上的 \co{seq}
	不是 \co{unsigned} 而是 \co{unsigned long}？
	毕竟，如果 \co{unsigned} 对 Linux 内核来说足够好，
	对每个人不应该也足够好吗？
}\QuickQuizAnswerE{
	完全不是。
	Linux 内核有许多特殊属性，使其可以忽略以下事件序列：
	\begin{enumerate}
	\item	线程~0 执行 \co{read_seqbegin()}，在
		\clnrefr{ln:defer:seqlock:impl:read_seqbegin:fetch}
		中获取 \co{->seq}，
		注意到该值为偶数，因此返回给调用者。
	\item	线程~0 开始执行其读侧临界区，
		但随后被抢占了很长时间。
	\item	其他线程反复调用 \co{write_seqlock()} 和
		\co{write_sequnlock()}，直到 \co{->seq} 的值
		溢出回到线程~0 获取的值。
	\item	线程~0 恢复执行，以不一致的数据完成
		其读侧临界区。
	\item	线程~0 调用 \co{read_seqretry()}，
		它错误地得出结论认为线程~0 看到了
		sequence lock 保护的数据的一致视图。
	\end{enumerate}

	Linux 内核对很少更新的东西使用 sequence locking，
	时间信息就是一个例子。
	此信息每毫秒最多更新一次，
	因此需要七周时间才能溢出计数器。
	如果一个内核线程被抢占了七周，
	Linux 内核的软锁检测代码会在整个期间
	每两分钟发出警告。

	相比之下，使用 64 位计数器，
	即使每\emph{纳秒}更新一次，
	也需要超过五个世纪才能溢出。
	因此，此实现使用在 64 位系统上为 64 位的类型作为 \co{->seq}。
}\QuickQuizEndE
}

\begin{listing}
\input{CodeSamples/defer/route_seqlock@lookup.fcv}
\caption{Sequence-Locked Pre-BSD 路由表查找（有缺陷！！！）}
\label{lst:defer:Sequence-Locked Pre-BSD Routing Table Lookup}
\end{listing}

\begin{listing}
\input{CodeSamples/defer/route_seqlock@add_del.fcv}
\caption{Sequence-Locked Pre-BSD 路由表添加\slash 删除（有缺陷！！！）}
\label{lst:defer:Sequence-Locked Pre-BSD Routing Table Add/Delete}
\end{listing}

那么当 sequence locking 应用于 Pre-BSD 路由表时会发生什么？
\Cref{lst:defer:Sequence-Locked Pre-BSD Routing Table Lookup}
展示了数据结构和 \co{route_lookup()}，
\cref{lst:defer:Sequence-Locked Pre-BSD Routing Table Add/Delete}
展示了 \co{route_add()} 和 \co{route_del()}（\path{route_seqlock.c}）。
此实现再次与前面几节中的对应物类似，因此只突出差异。

\begin{fcvref}[ln:defer:route_seqlock:lookup]
在\cref{lst:defer:Sequence-Locked Pre-BSD Routing Table Lookup}中，
\clnref{struct:re_freed}添加了 \co{->re_freed}，
在\clnref{lookup:chk_freed,lookup:abort}中被检查。
\Clnref{struct:sl}添加了 sequence lock，由 \co{route_lookup()} 使用
\end{fcvref}
\begin{fcvref}[ln:defer:route_seqlock:lookup:lookup]
在\clnref{r_sqbegin,r_sqretry1,r_sqretry2}上，
\clnref{goto_retry1,goto_retry2}分支回到
\clnref{retry}上的 \co{retry} 标签。
效果是重试任何与更新并发运行的查找。
\end{fcvref}

\begin{fcvref}[ln:defer:route_seqlock:add_del]
在\cref{lst:defer:Sequence-Locked Pre-BSD Routing Table Add/Delete}中，
\clnref{add:w_sqlock,add:w_squnlock,del:w_sqlock,%
del:w_squnlock1,del:w_squnlock2}
获取和释放 sequence lock，
而\clnref{add:clr_freed,del:set_freed}处理 \co{->re_freed}。
因此此实现相当直接。
\end{fcvref}

\begin{figure}
\centering
\resizebox{2.5in}{!}{\includegraphics{CodeSamples/defer/data/hps.2019.12.17a/perf-seqlock}}
\caption{受 Sequence Locking 保护的 Pre-BSD 路由表}
\label{fig:defer:Pre-BSD Routing Table Protected by Sequence Locking}
\end{figure}

它在只读工作负载上也表现更好，
如\cref{fig:defer:Pre-BSD Routing Table Protected by Sequence Locking}所示，
尽管其性能仍然远非理想。
更糟糕的是，它遭受释放后使用失败。
问题在于读者可能在 \co{read_seqretry()} 有机会警告并发更新之前，
由于访问已释放的结构而遇到段错误。

\QuickQuiz{
	这个错误能修复吗？
	换句话说，你能使用 sequence locks 作为\emph{唯一}
	的同步机制来保护支持并发添加、删除和查找的链表吗？
}\QuickQuizAnswer{
	一种微不足道的方式是用 \co{write_seqlock()} 和
	\co{write_sequnlock()} 包围所有访问，包括只读访问。
	当然，这个方案也禁止了所有读侧并行性，
	导致大量锁竞争，而且可以用简单的锁来同样轻松地实现。

	如果你确实想出了一种使用 \co{read_seqbegin()} 和
	\co{read_seqretry()} 来保护读侧访问的方案，
	请确保正确处理以下事件序列：

	\begin{enumerate}
	\item	CPU~0 正在遍历链表，并获取了指向列表元素~A 的指针。
	\item	CPU~1 从列表中移除元素~A 并释放它。
	\item	CPU~2 分配了一个不相关的数据结构，
		获得了之前被元素~A 占用的内存。
		在这个不相关的数据结构中，
		之前用于元素~A 的 \co{->next} 指针的内存
		现在被一个浮点数占用。
	\item	CPU~0 获取之前是元素~A 的 \co{->next} 指针的内容，
		得到随机位，因此遇到段错误。
	\end{enumerate}

	防止这类问题的一种方法需要使用``type-safe memory''，
	这将在\cref{sec:defer:Type-Safe Memory}中讨论。
	使用\cref{sec:defer:Hazard Pointers}中讨论的
	hazard pointers 也可以实现大致类似的解决方案。
	但无论哪种情况，你都将在 sequence locks 之外使用
	其他同步机制！
}\QuickQuizEnd

如\cpageref{sec:defer:Mysteries sequence locking}所暗示的，
sequence lock 的读侧和写侧临界区都可以被视为事务，
因此 sequence locking 可以被视为事务内存的一种受限形式，
这将在\cref{sec:future:Transactional Memory}中讨论。
Sequence locking 的局限性是：
\begin{enumerate*}[(1)]
\item Sequence locking 限制更新，
\item Sequence locking 不允许遍历可能被更新者释放的对象的指针。
\end{enumerate*}
这些局限性当然可以被事务内存克服，
但也可以通过将其他同步原语与 sequence locking 结合来克服。

Sequence locks 允许写者延迟读者，但反之则不然。
这可能导致\IX{unfairness}甚至在写密集型工作负载中的\IX{starvation}。\footnote{
	Dmitry Vyukov 描述了一种减少（但遗憾的是不能消除）
	读者饥饿的方法：
	\url{http://www.1024cores.net/home/lock-free-algorithms/reader-writer-problem/improved-lock-free-seqlock}。}
另一方面，在没有写者的情况下，sequence-lock 读者相当快且线性可扩展。
人们自然想要两全其美：
快速的读者且没有读侧失败的可能性，更不用说饥饿了。
此外，克服 sequence locking 在指针方面的局限性也会很好。
下一节将介绍一种恰好具有这些属性的同步机制。
